\section{Benchmark Circuits}\label{sec:algorithms}

All circuits are produced by one generator with fixed seeds; their gate counts and routing are generated from the
compiler counters and reported in Table~\ref{tab:routing}.

\paragraph{Quantum Fourier transform.} The $n$ qubit QFT consists of $n$ Hadamard gates, $n(n-1)/2$ controlled phase
gates $\mathrm{CP}(\pi/2^{j-i})$ and $\lfloor n/2\rfloor$ final SWAP gates, a total of $n(n+1)/2 + \lfloor n/2\rfloor$
gates~\cite{nielsen2010quantum}. In the baseline design with the two top qubits as partition qubits, the routing of
the QFT at $n=12$ is 59 local, 21 diagonal and 4 exchange gates (the two Hadamard gates on the partition qubits and the
two SWAP gates that cross the partition boundary), counted gate by gate on the circuit before lowering. On the engine
SWAP gates are relabellings and the counts of Table~\ref{tab:routing} apply. The QFT is also run on a Haar random input state and as a round trip (a random
product state, QFT, inverse QFT), whose output equals its input.

\paragraph{GHZ and Bell states.} The GHZ state is prepared by the textbook chain (H on qubit 0 and CX$(i,i+1)$), by a
fan out from the top qubit, and by the logarithmic depth construction of Cruz et al.~\cite{cruz2019efficient};
the Bell state is H on qubit 0 followed by CX from qubit 0 to qubit $n-1$.

\paragraph{Grover search.} For one marked state among $N = 2^n$ and $\theta = \arcsin(2^{-n/2})$, $k$ iterations give
the success probability $P(k) = \sin^2\bigl((2k+1)\theta\bigr)$. We use the standard iteration count
$k_{\mathrm{opt}} = \lfloor \pi/(4\theta) \rfloor$, which lies next to the first maximum of $P(k)$ and guarantees
$P(k_{\mathrm{opt}}) \ge 1 - 2^{-n}$~\cite{boyer1998tight,grover1996fast}; because $P(k)$ is periodic in $k$, larger
iteration counts can come closer to one, at several times the cost. All Grover runs use $k_{\mathrm{opt}}$, and we
report the measured target probability against $P(k_{\mathrm{opt}})$; for $n=4$ this probability is 0.9613. The oracle
and diffuser use one multi controlled $X$ gate each, which the engine executes as a single record.

\paragraph{Phase estimation (Shor).} The order finding component of Shor's algorithm~\cite{shor1997polynomial} is
run as quantum phase estimation of the modular multiplication by $a = 7$ modulo $N = 15$ (order $r = 4$) with a
counting register of $t$ qubits and four work qubits; the counting register must show four peaks of probability $1/4$
at multiples of $2^t/4$. To scale the phase estimation to the full width we also run it on the phase gate $P(2\pi\phi)$
with $\phi = 5/16$, whose $t$ qubit counting register must return $2^t\phi$ with probability one. A larger order
finding instance (for example $N = 21$) needs modular multipliers that we did not implement; the phase estimation on
the phase gate covers the same circuit structure at every width.

\paragraph{Random and worst case circuits.} Random circuits have layers of random U3 gates on every qubit and CX or
CZ gates on random disjoint pairs. The worst case for the distributed design applies H and RX gates to the top qubits
in every layer, so that every layer contains non diagonal gates on partition qubits.
